\item \textbf{¿Por qué no es posible la siguiente situación?}
Considere el sistema lente-espejo de la figura P35.55. La lente convergente tiene $f_L = 0.200\text{ m}$ y el espejo cóncavo tiene $f_M = 0.500\text{ m}$. Están separados una distancia $d = 1.30\text{ m}$ y un objeto se ubica a $p = 0.300\text{ m}$ de la lente. Al mover una pantalla a diferentes posiciones a la izquierda de la lente, un estudiante encuentra dos posiciones distintas de la pantalla que forman imágenes nítidas en ella: una por luz que viaja directamente a la izquierda a través de la lente y otra por luz que viaja a la derecha, se refleja en el espejo y luego pasa a la izquierda por la lente.